\section{Experiment Setup}
\label{sec:experiment_setup}

Section~\ref{sec:framework} defined the framework abstractly. This
section instantiates it experimentally: which SLMs sit at which stages
under the pre-registered DOE (Section~\ref{sec:experiment_setup_doe}),
which benchmark functions the framework is evaluated on
(Section~\ref{sec:experiment_setup_datasets}), how statistical inference
is drawn (Section~\ref{sec:experiment_setup_stats}), and how the
strict-AND policy is tested for robustness under threshold perturbation
(Section~\ref{sec:experiment_setup_sensitivity}). We also describe the
post-processing categorisation of every closure row into decision
reasons (Section~\ref{sec:experiment_setup_postproc}) and the capability
preservation protocol used to compare hetero and mono cells at
per-sample resolution (Section~\ref{sec:experiment_setup_capability}).
The human-evaluation protocol used to validate the framework's judge
signal against human raters is specified in
Section~\ref{sec:experiment_setup_human_eval}.

\subsection{Design of experiments}
\label{sec:experiment_setup_doe}

A cell is a stage-to-model assignment $a : \mathcal{S} \to \mathcal{M} \cup \{\bot\}$, where $\bot$ denotes an ablated stage. The DOE consists of 20 pre-registered cells across three strata:

\begin{equation}
  \mathcal{C}_{\textsf{DOE}} \;=\;
  \underbrace{\bigl\{a : a(\textsf{spec}) = a(\textsf{test}) = a(\textsf{code})\bigr\}}_{\text{Mono stratum (6 cells)}}
  \;\cup\;
  \underbrace{\bigl\{a : |\{a(s) : s \in \mathcal{S}\}| = 3\bigr\}}_{\text{Hetero stratum (11 cells)}}
  \;\cup\;
  \underbrace{\mathcal{C}_{\textsf{null}}}_{\text{Null stratum (3 cells)}}.
  \label{eq:doe}
\end{equation}

\begin{table}[t]
\centering
\small
\caption{Pre-registered 20-cell DOE over six open-weight SLMs and three pipeline stages ($6^3 = 216$ full design space, sampled purposively, not by fractional factorial). Three strata: 6 mono, 11 hetero, 3 null. Hypothesis is the rationale recorded at design time.}
\label{tab:doe_summary}
\begin{tabular}{@{}llllll@{}}
\toprule
Cell & Stratum & $L_{\textsf{spec}}$ & $L_{\textsf{test}}$ & $L_{\textsf{code}}$ & Hypothesis \\
\midrule
M1 & mono & llama3.2 & llama3.2 & llama3.2 & Small-dense self-consistency floor (3\,B Meta). \\
M2 & mono & phi4 & phi4 & phi4 & Mid-dense reasoning-tuned mono baseline (14\,B Microsoft). \\
M3 & mono & qwen3.6 & qwen3.6 & qwen3.6 & Dense mono baseline (27\,B Alibaba). \\
M4 & mono & gemma4 & gemma4 & gemma4 & MoE mono baseline (26\,B Google). \\
M5 & mono & mistral-s.\ 3.2 & mistral-s.\ 3.2 & mistral-s.\ 3.2 & Function-calling-tuned dense mono (24\,B Mistral). \\
M6 & mono & qwen3-coder & qwen3-coder & qwen3-coder & Code-specialised MoE self-consistency (30\,B Alibaba). \\
H1 & hetero & phi4 & qwen3-coder & qwen3-coder & Predicate-reasoner at spec plus coder at test/code. \\
H2 & hetero & qwen3.6 & phi4 & qwen3-coder & Dense spec, reasoner tests, coder synth. \\
H3 & hetero & gemma4 & mistral-s.\ 3.2 & qwen3-coder & Cross-family triple: Google $\to$ Mistral $\to$ Alibaba-coder. \\
H4 & hetero & qwen3-coder & qwen3-coder & llama3.2 & Cheap drafter at synth: does a 3\,B model suffice? \\
H5 & hetero & llama3.2 & qwen3-coder & qwen3-coder & Cheap drafter at spec: can a 3\,B docstring survive downstream? \\
H6 & hetero & phi4 & qwen3.6 & phi4 & Phi sandwich with a different-family middle stage. \\
H7 & hetero & qwen3-coder & qwen3.6 & phi4 & Reverse-capability gradient: strongest first, weakest last. \\
H8 & hetero & gemma4 & qwen3-coder & mistral-s.\ 3.2 & MoE $\to$ MoE $\to$ dense: does architecture matching matter? \\
H9 & hetero & qwen3-coder & qwen3.6 & qwen3-coder & Strong-sandwich: MoE bookends dense for synthesis stability. \\
H10 & hetero & mistral-s.\ 3.2 & phi4 & qwen3-coder & Best-of-each-family-stage: Mistral spec + Phi test + Qwen-coder synth. \\
H11 & hetero & phi4 & gemma4 & qwen3-coder & Phi spec + Gemma test + Qwen-coder synth (alt family triple). \\
N1 & null & llama3.2 & llama3.2 & llama3.2 & Word-shuffled input: expected metric $\approx 0$ if metric is sound. \\
N2 & null & (skip) & qwen3-coder & qwen3-coder & Spec-stage ablation: quantifies $L_{\textsf{spec}}$ contribution. \\
N3 & null & qwen3-coder & (skip) & qwen3-coder & Test-stage ablation: quantifies $L_{\textsf{test}}$ contribution. \\
\bottomrule
\end{tabular}
\end{table}


\begin{table}[t]
\centering
\caption{Small Language Models in the Chapter-3 pipeline. All models are open-weight and served via Ollama; total parameter count is $<30$ B per the §1.1 SLM definition.}
\label{tab:model_lineup}
\begin{tabular}{@{}lccccc@{}}
\toprule
Role & Ollama tag & Family & Size (B) & Architecture & Year \\
\midrule
Pipeline & llama3.2:3b & Meta & 3.0 & dense & 2024 \\
Pipeline & phi4:14b & Microsoft & 14.0 & dense & 2024 \\
Pipeline & qwen3.6:27b & Alibaba & 27.0 & dense & 2026 \\
Pipeline & gemma4:26b & Google & 26.0 & MoE & 2026 \\
Pipeline & mistral-small3.2:24b & Mistral & 24.0 & dense & 2025 \\
Pipeline & qwen3-coder:30b & Alibaba-coder & 30.0 & MoE & 2025 \\
Judge & deepseek-r1:14b & DeepSeek & 14.0 & dense & 2025 \\
\bottomrule
\end{tabular}
\end{table}


Table~\ref{tab:doe_summary} presents the complete DOE with the pre-registered hypothesis for each cell. The full 5-model $\times$ 3-stage design space contains $6^3 = 216$ cells; we sample 20 as a hypothesis-driven fractional factorial. The mono stratum establishes each SLM's self-consistency floor as a baseline. The hetero stratum assigns three different SLMs across the three stages based on per-stage strengths derived from the second companion paper's per-operator capability map (phi-4 dominates predicate reasoning; qwen-coder dominates comparison operators). The null stratum comprises three control cells: N1 (word-shuffled input; expected metric $\approx 0$ if the closure metric is sound), N2 (spec-stage ablated; $L_{\textsf{spec}} = \bot$), and N3 (test-stage ablated; $L_{\textsf{test}} = \bot$).

The six pipeline SLMs (Table~\ref{tab:model_lineup}) are all open-weight and under 30~B parameters, ensuring the entire experiment fits within a single Colab~Pro+ A100/T4 monthly allowance. Five distinct model families are represented: Meta with Llama-3.2-3B~\citep{Llama32}, Microsoft with Phi-4-14B~\citep{Phi4}, Alibaba with Qwen3.6-27B~\citep{Qwen36} and Qwen3-Coder-30B~\citep{QwenCoder}, Google with Gemma-4-26B~\citep{Gemma4}, and Mistral with Mistral-Small-3.2-24B~\citep{MistralSmall}. DeepSeek-R1:14B~\citep{DeepSeekR1} serves as the family-disjoint judge $L_J$.

\subsection{Datasets}
\label{sec:experiment_setup_datasets}

The primary dataset (\emph{core}) is a stratified 150-function sample from HumanEval (45~functions) and MBPP (105~functions), drawn with seed~42. This preserves the MBPP/HumanEval ratio of the 300-function second-paper sample and enables direct cross-paper comparisons via shared per-function difficulty draws.

Two held-out subsets were prepared for contamination sensitivity analysis: a 25-problem LiveCodeBench~\citep{Jain2025LiveCodeBench} subset with publication date $\geq$ 2024-12-01 (post-training-cutoff for all evaluated SLMs) and a 50-problem HumanEval-Mutated subset obtained by applying function-rename, docstring-paraphrase, and parameter-permutation transformations to a subset of HumanEval. Related benchmark alternatives include BigCodeBench~\citep{Zhuo2025BigCodeBench} and CRUXEval~\citep{Gu2025CruxEval}, both released in 2025 and offering complementary decontamination properties. The held-out subsets were prepared but not fully swept in this study; contamination sensitivity results on them will be reported in an addendum.

\subsection{Statistical methodology}
\label{sec:experiment_setup_stats}

\subsubsection{Analysis of variance}

For each closure metric $Y$ (kill rate, pass rate, BERTScore stacked across paths), we fit a Type-III analysis-of-variance model

\begin{equation}
  Y_{ij} \;=\; \mu \;+\; \alpha_i \;+\; \beta_j \;+\; \varepsilon_{ij},
  \quad \varepsilon_{ij} \sim \mathcal{N}(0, \sigma^2),
  \label{eq:anova}
\end{equation}

where $\alpha_i$ is the fixed effect of cell $i \in \mathcal{C}_{\textsf{DOE}}$, $\beta_j$ is the fixed effect of sample $j$, and $\varepsilon_{ij}$ is residual error. To further test whether the three closure paths carry non-redundant information about cell performance, we extend Eq.~\ref{eq:anova} to include a path factor and its interaction with cell:

\begin{equation}
  Y_{ijk} \;=\; \mu \;+\; \alpha_i \;+\; \eta_k \;+\; (\alpha \cdot \eta)_{ik} \;+\; \beta_j \;+\; \varepsilon_{ijk},
  \label{eq:anova_interact}
\end{equation}

where $\eta_k$ is the fixed effect of path $k \in \{1, 2, 3\}$ and $(\alpha \cdot \eta)_{ik}$ is the interaction. A significant interaction validates that closure paths are not merely substitutable measurements of the same underlying quantity. Because \texttt{scipy.stats}--\texttt{statsmodels} version drift blocked the standard OLS route on the analysis host, we implement Eq.~\ref{eq:anova_interact} directly using QR-based least squares and Type-I sequential sums of squares in pure NumPy.

\subsubsection{Pairwise cell comparisons}

Pairwise mean differences between cells $i$ and $k$ are tested via Tukey's honest significant difference:

\begin{equation}
  q_{ik}
  \;=\;
  \frac{\bar{Y}_i - \bar{Y}_k}{\sqrt{MS_{\textsf{within}} / n_h}},
  \label{eq:tukey}
\end{equation}

where $n_h$ is the harmonic mean of per-cell sample sizes (mono cells have $n = 150$; hetero and null cells have $n = 75$). A pair is significant at family-wise $\alpha = 0.05$ after Bonferroni correction over the $20 \times 19 / 2 = 190$ pairwise comparisons.

\subsubsection{Judge--metric correlation}

The Pearson correlation between the automated metric $Y$ and the judge rating $J$ per path,

\begin{equation}
  r_p \;=\; \frac{\textstyle\sum_{r} (Y_r - \bar{Y})(J_r - \bar{J})}{\sqrt{\textstyle\sum_r (Y_r - \bar{Y})^2 \cdot \textstyle\sum_r (J_r - \bar{J})^2}},
  \label{eq:corr}
\end{equation}

quantifies the extent to which each path's automated metric agrees with the judge SLM's semantic-equivalence rating. Rows with parse failures ($J = -1$) are excluded.

\subsubsection{Mixed-effects logit for per-stage decomposition}

To isolate which stage assignment drives closure, we treat each cell as the product of three stage choices and fit a generalised linear mixed model. Let $\text{closure}_{ij} \in \{0, 1\}$ denote whether function $j$ closed under cell $i$. We model

\begin{equation}
  \text{logit}\bigl(\Pr[\text{closure}_{ij} = 1]\bigr)
  \;=\;
  \beta_0
  + \beta_1 \cdot a_i(\textsf{spec})
  + \beta_2 \cdot a_i(\textsf{test})
  + \beta_3 \cdot a_i(\textsf{code})
  + \gamma_j,
  \label{eq:mixedlogit}
\end{equation}

where $\gamma_j \sim \mathcal{N}(0, \sigma_{\textsf{sample}}^2)$ is a random intercept for sample. Estimates $\hat{\beta}_1, \hat{\beta}_2, \hat{\beta}_3$ quantify each stage's contribution to closure odds.

\subsubsection{Judge--human agreement}

For each binarised rating (judge threshold $J \geq \rho$ versus $J < \rho$), Cohen's $\kappa$ is

\begin{equation}
  \kappa
  \;=\;
  \frac{p_o - p_e}{1 - p_e},
  \label{eq:kappa}
\end{equation}

where $p_o$ is the observed agreement proportion between the judge SLM and a human annotator and $p_e$ is the chance-expected proportion under independence. Values follow Landis and Koch's bands: $\kappa > 0.6$ substantial, $\kappa > 0.8$ almost perfect. Effect-size conventions throughout the paper follow \citet{Cohen2013StatPower}. For the multi-annotator setting (five human raters + one judge in the present study), we additionally compute Krippendorff's $\alpha$~\citep{Marzi2024KrippendorffAlpha} for ordinal ratings $r_{ij}$ from rater $i$ on item $j$:

\begin{equation}
  \alpha
  \;=\;
  1 \;-\; \frac{D_o}{D_e},
  \quad
  D_o \;=\; \tfrac{1}{n}\sum_{j} \tfrac{1}{\binom{m_j}{2}}
            \sum_{i_1 < i_2} \delta(r_{i_1 j}, r_{i_2 j})^2,
  \label{eq:alpha}
\end{equation}

where $\delta$ is the ordinal-difference metric and $D_e$ is the expected disagreement under random pairing. Values $\alpha > 0.667$ are considered acceptable for inferential use. Both $\kappa$ and $\alpha$ are reported in the human-evaluation addendum (Section~\ref{sec:results_human_eval}); the headline numbers on the 60-pair five-annotator sample are $\alpha_{\textsf{ord}} = 0.193$ across the three human annotators and $\kappa_{\textsf{quad}} = 0.074$ between the judge SLM and the majority human rating.

\subsubsection{Sensitivity analysis}
\label{sec:experiment_setup_sensitivity}

To defend against threshold-choice attacks, closure rates are computed and reported at multiple $(\tau, \rho)$ combinations. For the mutation kill rate on Path~1, $\tau \in \{0, 0.5, 0.8\}$; for the judge threshold, $\rho \in \{2, 3, 4\}$. Ordinal stability of the per-cell ranking across thresholds is quantified using Kendall's $\tau_{\textsf{rank}}$ between the reference threshold $(\tau=0, \rho=3)$ and each alternative.

\subsubsection{Post-processing categorisation}
\label{sec:experiment_setup_postproc}

Every closure row is post-hoc labelled with the Algorithm~\ref{alg:validity} decision reason and the Algorithm~\ref{alg:test_filter} filter reason. The distribution of decision reasons per path (Section~\ref{sec:discussion_rq2}) is our primary judge--metric disagreement diagnostic; the distribution of filter reasons per (cell, path) triple (Section~\ref{sec:discussion_rq3}) is our primary diagnostic for the phi-4-in-test-stage attribution.

\subsubsection{Capability preservation}
\label{sec:experiment_setup_capability}

For each path, we identify the strongest mono baseline (M4 for Paths~1 and~3; M3 for Path~2) and restrict analysis to samples on which that baseline achieves valid closure. For every other cell, the \emph{capability-preservation rate} is the fraction of those same samples on which the cell also achieves valid closure. Values near~$1$ indicate that the cell does not sacrifice mono strengths; substantially lower values expose cells where specialisation costs capability on samples the strongest mono handles cleanly.

\subsection{Human-evaluation protocol}
\label{sec:experiment_setup_human_eval}

To validate the framework's judge signal against human raters, we
pre-registered a 60-pair three-annotator study drawn from the completed
sweep. Sampling is stratified into three 20-pair buckets: (a) the
\emph{frontier} bucket, drawn uniformly from cells with mid-range
closure rates and rows where the judge and metric agree on validity;
(b) the \emph{agreement} bucket, rows where both signals mark the row
valid or both mark invalid; and (c) the \emph{disputed} bucket, rows
where the judge and metric disagree. All 60 pairs are drawn from the
core-sweep results TSV; the pair sample is committed to the replication
package (\texttt{results/human\_eval\_pairs\_60.tsv}) and independent of
the annotators.

Annotators rate each pair on the same $0$--$4$ rubric used by the judge
SLM: 4~=~identical, 3~=~equivalent (same observable behaviour),
2~=~approximately equivalent, 1~=~clearly different, 0~=~unrelated
(Section~\ref{sec:framework}). Each annotator sees the pairs in a
per-annotator deterministic order (SHA-256 seed on the annotator
identifier) and is blinded to cell identity, closure path, judge
rating, and bucket assignment. Ratings and one-sentence written
justifications are captured through a Streamlit annotation app; the
app auto-saves after each rating and is fully resumable across
sessions.

The pre-registered analysis plan computes: (i) Krippendorff's ordinal
$\alpha$ (Eq.~\ref{eq:alpha}) across all human annotators; (ii)
pairwise linear- and quadratic-weighted Cohen's $\kappa$
(Eq.~\ref{eq:kappa}) for each pair of raters; (iii) Cohen's $\kappa$
between the judge SLM and the majority-vote human rating (ties broken
toward the higher category); (iv) coarse-resolution within-1 agreement
rates (fraction of pairs where two ratings differ by at most one
category); and (v) an FPR/FNR ablation of metric-only vs.\ judge-only
vs.\ strict-AND decision policies against the majority-human ground
truth. Results are reported in Section~\ref{sec:results_human_eval}.


\emph{End of Section~\ref{sec:experiment_setup}.} Section~\ref{sec:results}
reports the empirical findings anchored to
Eqs.~\ref{eq:killrate}--\ref{eq:corr} and Table~\ref{tab:doe_summary},
organised by research question.


